\documentclass[a4paper,12pt,titlepage]{article}

\usepackage{enumitem,xcolor,multicol,amssymb,caption}
\usepackage{colortbl,array,natbib,graphicx,amsmath}
%\usepackage{subfig}
\usepackage[margin=2.0cm]{geometry}
\usepackage{natbib}
\usepackage{url}
\usepackage{threeparttable}
\usepackage{longtable}
\usepackage{rotating,pdflscape}
\usepackage{subcaption}
\usepackage[toc,page]{appendix}


\captionsetup[figure]{font=small,labelfont=small}
\setcounter{tocdepth}{2}
\setcounter{secnumdepth}{2}

\newcolumntype{C}[1]{>{\centering\let\newline\\\arraybackslash\hspace{0pt}}m{#1}}

\newcommand{\mnras}{MNRAS}
\newcommand{\apjl}{ApJ}
\newcommand{\apj}{ApJ}
\newcommand{\apjs}{ApJS}
\newcommand{\aj}{AJ}
\newcommand{\nat}{Nature}
\newcommand{\aap}{A\and A}
\newcommand{\iaucirc}{IAU Circular}
\newcommand{\pasj}{PASJ}
\newcommand{\pasa}{PASA}
\newcommand{\apss}{Ap\and SS}
\newcommand{\araa}{ARA\and A}
\newcommand{\pasp}{PASP}
\newcommand{\na}{NA}
\newcommand{\bain}{BAIN}
\newcommand{\procspie}{Proceedings of the SPIE}
\newcommand{\nar}{NAR}
\setcounter{secnumdepth}{3}

\title{An Honours Thesis in Curtin Astronomy}

\author{Name Student Number\\Supervisor: Supervisor Name 1\\Co-Supervisor: Supervisor Name 2} 

\begin{document}

\maketitle



\clearpage
\thispagestyle{empty}

\begin{minipage}[H]{\textwidth}
\begin{center}
 {\LARGE \textbf{Thesis Title: this is an alternative header page}}
 \\[1.5cm]
 {\Large 
 Name}\\[0.07cm]
 {\large
 (Curtin ID: XXXXX})\\[0.5cm]
 {\Large
 DATE \\[2cm]}
 {\Large
 Curtin University\\
 %(Department)\\
 International Centre for Radio Astronomy Research\\[0.5cm]
 \begin{tabular}{rl}
 Supervisor:& Name 1\\
 Co-supervisor:& Name 1
 \end{tabular}}
 \end{center}
 \end{minipage}


\newpage
\tableofcontents
\newpage
\begin{abstract}

An overview of my work.


\end{abstract}

\newpage

\setcounter{page}{1}
\section{Introduction}\label{sec:intro}

%The majority of studies surrounding BHs rely on examining how the presence of a BH influences the environment in which it lives.

An introduction

\section{My methodology}


\subsection{A subsection of my methodology}

A figure with label \ref{ionisation}.

\begin{figure}
    \centering
    \includegraphics[width=0.5\columnwidth]{ionisation}
    \caption{Figure taken from \cite{1989agna.book.....O}.}
    \label{ionisation}
\end{figure}

A table with label \ref{optical}.

\begin{table}
\centering
\caption{Full galaxy sample and properties.}
\begin{threeparttable}

\begin{tabular}{l c r c c r c c}
\hline\hline
I.D\tnote{1} & Name\tnote{2} & D\tnote{3} & \textit{Z}\tnote{4} & $M_*$\tnote{5} & SFR\tnote{6} &  X-ray\tnote{7} & Alternate Name\\
  &  (SDSS) &   &  &   &   &  Sources &   \\
\hline
72 \tnote{8}& J024240.70$-$000047.9  &  18.6  &  8.93$\pm$0.25& 7.39  & -0.27  &  $>2$ & M077\\ %  02:42:39.86 & -00:00:58.64 &
21 & J083743.48$+$513830.2  &  10.8  & 7.95$\pm$0.08 & 6.04  & -1.84  &   0 & Mrk 0094 \\ % 08:37:43.48 & 51:38:30.26 &
101 & J101624.51+375445.9  &  17.2 & 7.80$\pm$0.05  & 6.58  & -1.50  &  0 & --- \\ % 10:16:24.53 & 37:54:45.97 &
62 & J103410.15+580349.1  & 33.0  &  7.78$\pm$0.06 &  7.00 & -1.21  &  2 & Mrk 1434 \\ % 10:34:10.15 & 58:03:49.06 &
124 & J104653.98+134645.7  &  46.9&  7.92$\pm$0.09 & 6.81  &  -1.15 &  0 & --- \\ % 10:46:53.99 & 13:46:45.77 &
59 & J105308.48+501704.9  & 19.2 & 7.94$\pm$0.08 &  6.04 & -2.12  &  1 & ---\\ % 10:53:10.82 & 50:16:53.21 & 
178 & J111746.29+174424.6 &  21.8 &  7.90$\pm$0.09 &  6.62 & -1.44  & 2 & ---\\ % 11:17:46.30 & 17:44:24.69 &
149 & J114107.48+322537.2  & 26.6  & 7.79$\pm$0.05  & 6.04  & -1.80  &  1 & KUG 1138+327\\ % 11:41:07.49 & 32:25:37.22 &
8 & J115237.67$-$022806.3  & 15.6  & 7.84$\pm$0.09  & 6.71  & -1.45  & 2 & --- \\ % 11:52:37.68	& -02:28:06.39 &
147 & J121749.30+375155.5  & 2.7  &  7.92$\pm$0.09 & 6.02  & -2.59  &  0 & --- \\
\hline
\hline

\end{tabular}
\end{threeparttable}
\label{optical}
\end{table}

An example equation

\begin{equation}
  \psi(M_{\odot} \text{ yr}^{-1})=
  \begin{cases}
    5.52\times10^{-22}L_{1.4\text{GHz}}, & L>L_c\\
    \frac{5.52\times10^{-22}}{0.1+0.9(L/L_c)^{0.3}}L_{1.4\text{GHz}}, & L\leq L_c
  \end{cases}
  \label{eq:sfr-Lr}
\end{equation}


\section{Conclusion and Future Work}\label{sec:conlusion}

\begin{appendix}
\section{Appendix}\label{sec:app}

This is where we put appended information.

\end{appendix}



\clearpage

\setlength{\bibsep}{0pt}
\bibliography{references}
\bibliographystyle{apj}
\end{document}




